\documentclass[12pt,a4paper]{report}
\usepackage[margin=1in]{geometry}
\usepackage{hyperref}
\usepackage[utf8]{inputenc}
\usepackage[T1]{fontenc}
\usepackage{enumitem}
\usepackage{booktabs}
\usepackage{listings}
\usepackage{xcolor}
\usepackage{graphicx}
\usepackage{fancyhdr}
\usepackage{titlesec}
\usepackage{parskip}

% ── Header / Footer ───────────────────────────────────────────
\pagestyle{fancy}
\fancyhf{}
\fancyhead[L]{\small UCS503 - Software Engineering}
\fancyhead[R]{\small SyncDrive}
\fancyfoot[C]{\thepage}
\renewcommand{\headrulewidth}{0.4pt}

% ── Listings style ────────────────────────────────────────────
\lstset{
  basicstyle=\ttfamily\small,
  breaklines=true,
  frame=single,
  backgroundcolor=\color{gray!8},
  rulecolor=\color{gray!30},
  xleftmargin=4pt, xrightmargin=4pt,
  aboveskip=6pt, belowskip=4pt
}

% ── Bibliography ──────────────────────────────────────────────
\usepackage[backend=biber,style=alphabetic,citestyle=alphabetic]{biblatex}
\addbibresource{references.bib}

% ─────────────────────────────────────────────────────────────
\begin{document}

% ════════════════════════════════════════════════════════════
%  COVER PAGE
% ════════════════════════════════════════════════════════════
\begin{titlepage}
\centering
\vspace*{2cm}

{\large\textbf{UCS503 - Software Engineering}}

\vspace{1.5cm}

{\Huge\textbf{SyncDrive}}\\[0.4cm]
{\large Signal-Aware Navigation and Speed Optimization System}

\vspace{3cm}

\textbf{Submitted by:}\\[0.3cm]
1024030177\quad Tarun Mahindra\\
1024030178\quad Parth Puri\\
1024031077\quad Adhiraj Kapur\\[0.4cm]
Group: BE Third Year, CSE

\vspace{2.5cm}

\textbf{Submitted to:}\\[0.3cm]
Dr.\ Paramveer Sidhu

\vspace{2.5cm}

Mid-Semester Evaluation\\[0.2cm]
\textbf{Computer Science and Engineering Department}\\
Thapar Institute of Engineering and Technology, Patiala

\end{titlepage}

% ════════════════════════════════════════════════════════════
%  BLANK PAGE
% ════════════════════════════════════════════════════════════
\newpage\thispagestyle{empty}\mbox{}

% ════════════════════════════════════════════════════════════
%  TABLE OF CONTENTS
% ════════════════════════════════════════════════════════════
\tableofcontents
\newpage

% ════════════════════════════════════════════════════════════
%  CHAPTER 1 — INTRODUCTION
% ════════════════════════════════════════════════════════════
\chapter{Introduction}

\section{Background and Context}

Urban commuters in Indian cities spend a significant portion of their travel
time waiting at traffic signals. In a city like Chandigarh, whose grid-sector
layout concentrates traffic through a small number of major intersections, the
stop-and-go pattern is both predictable and avoidable. Standard navigation
applications --- Google Maps, Apple Maps, and their equivalents --- select
routes based on distance and estimated travel time, but they do not account for
how a driver's chosen speed interacts with the timing of upcoming traffic
signals. A driver following the fastest route at the speed limit may still stop
at every signal, losing more time to red-light waits than would be saved by
taking a longer but signal-friendly path at a slightly lower speed.

The concept of a \textit{green wave} --- coordinating signal timings along a
corridor so that a vehicle travelling at a fixed speed passes every light on
green --- is well established in traffic engineering. City authorities use it to
improve arterial throughput. SyncDrive inverts this concept for the individual
driver: rather than the city adjusting signals to match a speed, the driver
adjusts their speed and departure time to match the existing signal pattern.

\subsection{Problem Statement}

Existing navigation systems are \textit{signal-blind}: they select a road but
not the speed or departure time that would allow the driver to traverse that
road without stopping. This results in three compounding inefficiencies:

\begin{itemize}
  \item \textbf{Avoidable red-light stops:} a driver who arrives at a signal
    thirty seconds after it turned red waits up to a full cycle (60--90
    seconds) for no mechanical reason.
  \item \textbf{Stop-and-go fuel waste:} repeated acceleration from rest
    consumes significantly more fuel per kilometre than steady-speed driving.
  \item \textbf{No departure-time guidance:} apps tell a driver \textit{which
    road} to take, not \textit{when to leave} or \textit{how fast to drive} on
    it.
\end{itemize}

The core claim of SyncDrive is that a driver travelling at a moderate, steady
speed with zero red-light stops can arrive faster than one who rushes at the
maximum speed but accumulates waiting time at signals.

\section{Project Objectives}

\begin{enumerate}
  \item \textbf{Objective 1 --- Green wave optimization:} Given a source,
    destination, and desired arrival time, compute the departure time and
    steady driving speed that minimises total journey time (drive time plus
    signal wait time) across all signals on the route.

  \item \textbf{Objective 2 --- Quantified comparison:} Simulate a Normal GPS
    journey on the same route at the same departure time but at the city speed
    limit (50 km/h), compute the red-light wait time accumulated, and present
    the time saved by SyncDrive side by side.

  \item \textbf{Objective 3 --- Working browser prototype:} Deliver a
    single-file web application --- no install, no backend server, no API key
    --- that renders a dark interactive map of Chandigarh, allows the user to
    select from 7 named locations across 42 directed routes, and displays the
    full recommendation including departure time, speed, signal timeline, and
    turn-by-turn directions.
\end{enumerate}

All three objectives are specific, measurable, and verifiable within the
prototype: the optimizer produces a numeric departure time and speed; the
comparison produces numeric time values; the browser prototype either loads and
works or it does not.

% ════════════════════════════════════════════════════════════
%  CHAPTER 2 — METHODOLOGY AND DESIGN
% ════════════════════════════════════════════════════════════
\chapter{Methodology and Design}

\section{System Architecture}

The prototype is a single HTML file with embedded CSS and JavaScript. There is
no backend, no build step, and no external API dependency beyond map tile
fetching. The entire application loads and runs in the browser.

\subsection{High-Level Design}

The system has four logical layers:

\begin{enumerate}
  \item \textbf{Data Layer} --- Static JavaScript objects: \texttt{PLACES}
    (7 locations), \texttt{SIGNALS} (16 traffic signals), \texttt{ROUTES}
    (42 directed route objects).
  \item \textbf{Algorithm Layer} --- Two functions: \texttt{optimise()} (green
    wave optimizer) and \texttt{simNG()} (Normal GPS simulator).
  \item \textbf{Map Layer} --- Leaflet.js 1.9.4 rendering an interactive map
    with place labels, signal markers, route polylines, and a speed HUD.
  \item \textbf{UI Layer} --- Three-pane Apple Maps-inspired layout: nav rail,
    directions panel, map.
\end{enumerate}

\subsection{Technical Stack}

\begin{itemize}
  \item \textbf{Language:} Vanilla JavaScript (ES6+), HTML5, CSS3.
  \item \textbf{Map library:} Leaflet.js 1.9.4 (CDN).
  \item \textbf{Map tiles:} Esri World Dark Gray Canvas (free, no API key).
  \item \textbf{Font:} Inter via Google Fonts CDN.
  \item \textbf{Build tools:} None --- single \texttt{index.html} file.
  \item \textbf{Deployment:} GitHub Pages, served from repository root on the
    \texttt{master} branch.
  \item \textbf{Version control:} Git / GitHub
    (\url{https://github.com/ParthPuri/SyncDrive}).
\end{itemize}

\section{Implementation Details}

\subsection{Module 1: Green Wave Optimizer}

The optimizer searches over a grid of candidate (speed, departure offset) pairs
and returns the combination that minimises total journey time.

\textbf{Algorithm:}
\begin{lstlisting}
Speeds tested : [25, 28, 30, 32, 35, 38, 40, 42, 45] km/h
Offset window : 0 to 1500 seconds before latest departure
Step size     : 15 seconds

for each speed s:
  driveSec = (distKm / s) * 3600
  latestDep = arrivalTime - driveSec

  for each offset d in 0..1500 step 15:
    dep = latestDep - d
    elapsed = 0;  totalWait = 0;  greens = 0

    for each signal on route:
      phase = (dep + timeToSignal + elapsed) % cycle
      if green(signal, phase):
        greens++
      else:
        wait = nextGreenOffset(signal, phase)
        elapsed += wait;  totalWait += wait

    score = (driveSec + totalWait) - 0.5 * greens
    if score < best: save result
\end{lstlisting}

The \texttt{elapsed} accumulation is the critical correctness requirement: each
red-light stop pushes the driver's arrival at all downstream signals later,
potentially converting a green into a red. Without this propagation the signal
counts are wrong.

\textbf{Signal phase mathematics:}
\begin{itemize}
  \item \texttt{phase = (depTime + travelToSignal + elapsed) \% cycle}
  \item \texttt{isGreen = phase >= greenStart \&\& phase < (greenStart +
    greenDuration)}
  \item \texttt{waitIfRed = (phase < greenStart) ? (greenStart - phase) :
    (cycle - phase + greenStart)}
\end{itemize}

\subsection{Module 2: Normal GPS Simulation}

To produce a fair comparison, the Normal GPS simulator uses the \textit{same
departure time} as the SyncDrive result but drives at 50 km/h (city speed
limit). It hits signals at uncontrolled phases and accumulates real wait times
using the same \texttt{elapsed} propagation logic.

Using the same departure time is essential: if the Normal GPS simulation
departed later (i.e., at the last possible moment to reach the destination at
50 km/h), it would have less travel time and the comparison would be
misleading. Equal departure time means any difference in total time is
attributable solely to speed choice and signal timing --- not to when the
driver left.

% ════════════════════════════════════════════════════════════
%  CHAPTER 3 — RESULTS AND EVALUATION
% ════════════════════════════════════════════════════════════
\chapter{Results and Evaluation}

\section{Testing Methodology}

Testing was conducted computationally using the simulator built into the
prototype. All 42 directed routes were evaluated at a 09:00 arrival target.

\begin{itemize}
  \item \textbf{Unit testing:} Each signal phase function was verified manually
    against hand-computed values for known (departure, cycle, offset)
    combinations.
  \item \textbf{Integration testing:} The optimizer and simulator were run on
    all 42 routes and results were checked for monotonicity (longer routes
    should produce longer travel times) and correctness (GreenWave total should
    never exceed Normal GPS total for the same departure).
  \item \textbf{Boundary testing:} Departure times near midnight (0 seconds)
    were checked to confirm the phase modular arithmetic does not produce
    negative values.
  \item \textbf{Browser testing:} The application was loaded in Chrome and
    Safari and verified to render correctly with no console errors.
\end{itemize}

\section{Performance Metrics}

\begin{table}[h]
\centering
\begin{tabular}{lrrr}
\toprule
\textbf{Route} & \textbf{km} & \textbf{Normal GPS} & \textbf{SyncDrive} \\
\midrule
Sector 17 $\to$ Sector 22  & 1.4 & 3.8 min & 2.6 min \\
Sector 22 $\to$ Sector 35  & 2.2 & 6.0 min & 4.5 min \\
Sector 17 $\to$ PGI        & 3.2 & 5.8 min & 5.0 min \\
Sector 43 $\to$ Elante Mall & 2.0 & 4.1 min & 4.0 min \\
Sector 17 $\to$ Tribune Chowk & 5.1 & 8.2 min & 7.7 min \\
Sector 22 $\to$ PGI        & 4.1 & 7.9 min & 7.0 min \\
Sector 35 $\to$ Sector 43  & 2.8 & 5.9 min & 4.9 min \\
\bottomrule
\end{tabular}
\caption{Journey time comparison: Normal GPS vs SyncDrive (09:00 arrival target).
All figures are based on simulated signal timings.}
\end{table}

\begin{table}[h]
\centering
\begin{tabular}{lcc}
\toprule
\textbf{Metric} & \textbf{Target} & \textbf{Achieved} \\
\midrule
Routes covering all 7 locations & 42 & 42 \\
Signals modelled & $\geq$10 & 16 \\
SyncDrive beats or ties Normal GPS & 100\% of routes & 100\% \\
Application loads without server & Yes & Yes \\
No API key required & Yes & Yes \\
\bottomrule
\end{tabular}
\caption{Prototype objective achievement summary.}
\end{table}

\section{Analysis}

\textbf{Did we meet our objectives?} Yes. All three objectives from Chapter 1
are fully met. The optimizer produces a numeric departure time and speed for
every route pair. The comparison panel displays the time breakdown
side-by-side. The single-file prototype loads in any browser with no
installation.

\textbf{Where did we exceed expectations?} The original proposal described a
single test corridor. The prototype covers 7 locations, 16 signals, and 42
directed routes --- significantly broader than scoped. The UI additionally
renders a per-signal timeline, turn-by-turn directions, and a speed HUD on
the map.

\textbf{Challenges encountered and resolved:}

\begin{itemize}
  \item \textbf{Tile provider API keys:} Both Carto Voyager and Stadia Maps
    began requiring API keys after project start. Resolved by switching to
    Esri ArcGIS Dark Gray Canvas, which remains freely accessible without
    authentication.
  \item \textbf{Wrong optimization objective:} The first version of the
    optimizer maximised green signal count, which ignored total journey time.
    A route with 4 greens at 25 km/h could score higher than one with 3 greens
    at 45 km/h even if the latter was faster overall. Fixed by changing the
    objective to minimise $(t_{\text{drive}} + t_{\text{wait}})$.
  \item \textbf{Missing elapsed propagation:} Early builds did not accumulate
    wait time through the signal loop, causing incorrect downstream phase
    calculations. Fixed by threading the \texttt{elapsed} variable through the
    signal simulation loop.
  \item \textbf{Unfair comparison baseline:} The Normal GPS simulation
    initially used the arrival time as its departure reference, giving it extra
    time. Fixed to use the same departure time as SyncDrive.
\end{itemize}

% ════════════════════════════════════════════════════════════
%  CHAPTER 4 — CONCLUSION
% ════════════════════════════════════════════════════════════
\chapter{Conclusion}

\section{Summary of Work}

SyncDrive demonstrates that a driver can arrive at a destination faster by
driving at a lower, steady speed timed to a signal wave, compared to rushing
at the city speed limit and accumulating red-light wait time. The prototype
delivers a working single-file web application covering Chandigarh's key
corridors, with an interactive dark map, 42 routes, 16 simulated signals, a
green wave optimizer, a Normal GPS comparison, and an Apple Maps-inspired user
interface --- all without a backend server or API key.

\section{Key Findings}

\begin{itemize}
  \item \textbf{Finding 1 --- Slower speed, faster arrival:} On short routes
    (1--2 km), driving at 30--35 km/h with zero stops consistently beats
    50 km/h with 2--3 red-light stops by 0.5--1.5 minutes.
  \item \textbf{Finding 2 --- Elapsed propagation is non-negotiable:} The
    single most important correctness requirement in the simulator is
    accumulating wait time through the signal loop. Omitting it produces
    systematically optimistic green counts.
  \item \textbf{Finding 3 --- Fair comparison requires equal departure time:}
    Using the same departure time for both strategies isolates the effect of
    speed and signal timing, making the comparison scientifically valid and
    the time savings directly attributable to the green wave approach.
\end{itemize}

\section{Future Work and Recommendations}

\begin{enumerate}
  \item \textbf{Live position tracking:} Use the browser Geolocation API to
    track the driver's real-time position and continuously recalculate the
    recommended speed as the journey progresses \cite{OSM26}.
  \item \textbf{Departure countdown:} Add a ``Leave in X minutes'' live timer
    based on the current system clock.
  \item \textbf{Animated route playback:} Move a vehicle marker along the
    polyline at the recommended speed to make the green wave visually
    intuitive.
  \item \textbf{Real signal data:} Source signal cycle lengths from
    OpenStreetMap \texttt{highway=traffic\_signals} and
    \texttt{cycle\_time:*} tags as a richer, city-verified data source.
  \item \textbf{Peak-hour profiles:} Model different signal cycle lengths for
    morning (08:00--10:00) and evening (17:00--20:00) rush hours.
  \item \textbf{OSRM routing integration:} Replace hardcoded polylines with
    road-following routes from the OSRM open-source routing engine to improve
    geographic accuracy.
  \item \textbf{Emergency vehicle coordination:} Phase 2 feature from the
    original proposal --- adjust nearby drivers' speed recommendations when an
    emergency vehicle is active on the same corridor.
\end{enumerate}

\section{Final Remarks}

The prototype validates the central hypothesis of SyncDrive within a
controlled, simulated environment. The engineering challenge was not the
map rendering or the UI --- it was getting the signal phase simulation
mathematically correct and ensuring the comparison between strategies was
scientifically fair. Both issues were identified, debugged, and fixed during
the five weeks of development documented in the team journals. The foundation
is solid for the next phase of work.

% ════════════════════════════════════════════════════════════
%  BIBLIOGRAPHY
% ════════════════════════════════════════════════════════════
\printbibliography

\end{document}
